% !TeX root = ../papers/nc-p3.tex
\section{二変数の非分岐単元域全体の経路と四次中心}
\label{sec:rank-fifteen-bivariate-word-paths}
\begin{theorem}\label{thm:rank-fifteen-bivariate-word-paths}
標数3で $K,H$ を取り、
\[
 KH(K^2-1)\big((K^2-1)^2-4H^2\big)
 \big((K+1)^2-H^2\big)\big((K-1)^2-H^2\big)\ne0
\]
とする。36成分に規格化した窓の全始点で、長さ1から5までのKlein文字別経路次元は
\[
 (1,1,1,1),\quad (2,2,2,4),\quad (5,5,5,5),\quad
 (11,8,8,8),\quad(14,14,14,14)
\]
である。四次中心の文字別次元は $(5,0,0,0)$ であり、文字0の中心の同じ自由列が
開集合全体で基底を与える。この範囲には任意の幾何単元 $X\ne1$ と
$Y\ne0,(1-X)^2$ の表示が含まれる。
\end{theorem}
\begin{proof}
$F=\mathbb F_3(K,H)$ の二変数有理関数として独立に計算する。
$P=W^2$ を六根の二乗の積とし、前節の偶奇式を $\ell_p$ と書く。
\path{computations/verify_rank_fifteen_bivariate_word_paths.py} は全ての係数を
$\mathbb F_3[K,H]$ の互いに素な分子とモニック分母で持つ。

規格化した語 $W^{\ell_p(w)}w$ を、各長さ、始点、文字の商の基底に射影する。
先頭の矢を $a$ とし、末尾の語 $t$ の射影に現れる基底語を $s$ とする。
既存の $s$ の係数に掛ける補正は
\[
 P^{(\ell_p(at)-\ell_{p+1}(t)+\ell_{p+1}(s)-\ell_p(as))/2}
\]
である。$s,t$ の文字が等しいので分子は偶数である。
二次関係に固定した末尾を付ける際も、各項の $W$ 指数から行の共通偶奇を引いて
$P$ の整数冪を使う。これにより全関係行を構成し、主元行列を後退代入して
次の長さの語基底と射影を求める。各始点で表の次元が得られる。

文字0の四次中心方程式では未知数が $4\cdot11=44$ 個である。
各語の左右に矢を掛け、上の五次への射影で比較する。行共通の $W$ 冪を除くと
$F$ 上で階数39となる。ここまでに逆にした主元の分子の既約因子は、
$H,K-1,K+1$ と六根の二乗のモニック化の九つだけである。
他の三文字では交換子の左右で偶奇が違う場合がある。
各未知数を $u+Wv$ として $1,W$ 成分に分けると、各文字で $F$ 上64列となる。
上の九因子と $K$ に関して単元となる主元を優先して選べば、三つの行列とも階数64である。
従ってどの段階にも定理の開集合の内部で消える分母はない。
$W^2=P$ が特殊化で分解しても、$P\ne0$ なので二つの符号を合わせた行列が全列階数なら
両符号の核は零である。元の関係と全ての行操作が特殊化できるので、零行も保存される。

幾何単元表示で $K=1/v$, $H=c(K-v)$ とすると、$Y\ne0$ から $v\ne0,\pm1$、
従って $K,H,K\pm1$ は非零である。残る因子は六根の消失に対応し、これは分岐水準に限られる。
\end{proof}
